\chapter{Design Considerations and Standards}

\section{Governing Standards and Design Framework}

The design of the pump station and force main system for the Colebrooke Subdivision is governed by established regulatory standards and accepted engineering practices for wastewater conveyance systems. These standards provide the framework for ensuring that the system meets requirements related to hydraulic performance, operational reliability, and public health protection.

The primary regulatory guidance for this project is derived from the \textit{Guidance for the Design of Publicly Owned Wastewater Facilities} issued by the Mississippi Department of Environmental Quality (MDEQ) and the \textit{Recommended Standards for Wastewater Facilities} (Ten States Standards)~\cite{MDEQ_DesignGuidance_2025,TenStates_RecommendedStandards_2014}. These documents establish minimum criteria for the design of wastewater pumping systems, including requirements for hydraulic capacity, minimum flow velocities, pump redundancy, wet well operation, and system reliability.

In addition to regulatory standards, general design methodology follows established engineering practices for pump station design. According to the \textit{Pump Station Design Guidelines} published by Jensen Engineered Systems, the design process begins with defining the flow conditions and site constraints, followed by evaluation of system head requirements, development of the system curve, and selection of a pump that satisfies the hydraulic demands of the system~\cite{Jensen_PumpStationDesignGuidelines_2012}.

These steps emphasize the importance of accurately characterizing both the flowrate and the total dynamic head, as the intersection of the system curve and pump performance curve determines the operating point of the system. The methodology ensures that the selected pumping system operates efficiently and reliably under expected conditions.

Collectively, these standards and design principles establish the basis for the hydraulic analysis and system design presented in subsequent sections. The following sections further define the specific criteria and methods used to develop the final pump station and force main design.

\section{Hydraulic Design Criteria}

The hydraulic design of the pump station and force main system is governed by established criteria to ensure reliable conveyance of wastewater while preventing operational issues such as solids deposition, excessive wear, and energy inefficiency. These criteria define the acceptable range of flow conditions within the system and serve as the basis for subsequent design decisions.

A primary requirement for wastewater force mains is the maintenance of self-cleansing velocity to prevent the accumulation of solids within the pipeline. According to standard wastewater design guidelines, a minimum velocity of approximately 2~ft/s is required to ensure adequate transport of suspended solids and to minimize the risk of sedimentation and clogging~\cite{TenStates_RecommendedStandards_2014}. Maintaining this velocity is critical for preserving long-term operational efficiency and reducing maintenance requirements.

In addition to minimum velocity, upper limits on flow velocity are also considered to prevent excessive headloss and potential damage to system components. While specific limits may vary depending on material and system configuration, typical wastewater design practice maintains velocities within a practical range that balances solids transport and energy efficiency. Excessively high velocities can increase friction losses and energy consumption, while also contributing to pipe wear over time.

The hydraulic design is based on peak flow conditions, as defined in Section~2.3, to ensure that the system performs adequately under maximum loading scenarios. Designing for peak flow provides a conservative basis for system sizing and ensures that the conveyance system can accommodate fluctuations in wastewater generation.

These hydraulic criteria establish the performance requirements for the force main and pumping system. The selected pipe diameter and pump configuration must satisfy these conditions while maintaining efficient and reliable operation. The application of these criteria is demonstrated in the hydraulic analysis presented in later chapters.


\section{Headloss Calculation Method}

The hydraulic performance of the force main system is evaluated by determining the total dynamic head (TDH) required to convey wastewater from the pump station to the downstream municipal sewer connection. The TDH represents the total energy that must be supplied by the pump and is composed of static head, friction losses, and minor losses within the system.

The static head is defined by the elevation difference between the wastewater level in the wet well and the discharge point at the municipal sewer system. In addition to this elevation component, energy losses occur due to friction between the flowing fluid and the internal surfaces of the pipe, as well as through fittings, valves, and changes in flow direction.

For this study, friction losses in the force main are estimated using the Hazen--Williams equation, which is commonly applied in water and wastewater conveyance systems due to its simplicity and practical applicability. The Hazen--Williams method is widely accepted in engineering practice and is recommended for systems involving water or wastewater under typical operating conditions.

The general form of the Hazen--Williams equation relates headloss to flow rate, pipe diameter, and a material-dependent roughness coefficient. This method assumes turbulent flow conditions and is most appropriate for relatively simple piping systems such as a single force main conveying wastewater to a gravity discharge point.

Minor losses associated with fittings, valves, and transitions are accounted for using standard loss coefficients, which are applied to estimate additional headloss contributions within the system. These losses, although smaller than friction losses, are important for accurately determining the total energy requirement.

The total dynamic head is therefore expressed as the sum of static head, friction losses, and minor losses. This relationship forms the basis for developing the system curve, which is used in conjunction with pump performance curves to select an appropriate pumping system~\cite{Jensen_PumpStationDesignGuidelines_2012}.

The application of these methods to the Colebrooke Subdivision system is presented in the hydraulic analysis and design calculations in Chapter~5.



\section{Pump Station Design Criteria}

The design of the pump station must satisfy operational, hydraulic, and reliability requirements to ensure continuous and efficient wastewater conveyance. Wastewater lift stations are critical infrastructure components, and their performance directly affects system functionality and public health protection.

A fundamental requirement for pump station design is the provision of redundancy. In accordance with standard wastewater design guidelines, pump stations are typically configured with multiple pumps to ensure uninterrupted operation in the event of equipment failure or maintenance. A duplex pump configuration, consisting of two pumps operating in an alternating duty cycle, is commonly adopted for small to medium-sized systems~\cite{TenStates_RecommendedStandards_2014,MDEQ_DesignGuidance_2025}. This configuration allows one pump to operate under normal conditions while the second pump remains available as standby or operates during peak flow conditions.

Wet well design is also an essential component of pump station operation. The wet well must be sized to accommodate variations in inflow while minimizing excessive pump cycling. Proper control of water levels within the wet well ensures efficient pump operation and extends equipment life by preventing frequent start-stop cycles~\cite{TenStates_RecommendedStandards_2014}.

Pump selection must consider both the required flowrate and the total dynamic head of the system. The selected pump must operate within an efficient range while meeting the hydraulic demands defined by the system curve. Reliable operation under varying flow conditions is necessary to maintain consistent wastewater transport. The general methodology for pump selection, including the use of system curves and pump performance curves, follows established design practices outlined in pump station design guidelines~\cite{Jensen_PumpStationDesignGuidelines_2012}.

Additional considerations include ease of maintenance, accessibility of equipment, and operational safety. Submersible pump configurations are commonly used in wastewater applications due to their compact design, ability to operate under submerged conditions, and suitability for handling wastewater containing suspended solids~\cite{MDEQ_DesignGuidance_2025}.

These criteria establish the requirements that the final pump station design must satisfy. The application of these principles is presented in Chapter~5, where the pump configuration and operating conditions are evaluated based on the system requirements.


\section{Material Selection Criteria}

Material selection for the pump station and force main system must ensure durability and compatibility with wastewater service conditions. A primary requirement is resistance to corrosion and chemical degradation, as wastewater environments contain constituents that can deteriorate susceptible materials over time. In addition, materials must provide adequate structural performance to withstand external loads from soil, groundwater, and surface conditions while maintaining integrity and alignment of buried infrastructure~\cite{MDEQ_DesignGuidance_2025,TenStates_RecommendedStandards_2014}.

Hydraulic efficiency and constructability also influence material selection. Materials with smooth internal surfaces reduce friction losses and improve flow performance, which is important for maintaining efficient pumped operation. At the same time, selected materials must be practical for installation under site-specific conditions, including trenching and groundwater constraints, and must be compatible with standard fittings and system components. These considerations ensure that the final material selection achieves a balance between performance, reliability, and long-term cost-effectiveness~\cite{MDEQ_DesignGuidance_2025}.


\section{Operational and Safety Considerations}

The pump station and force main system must be designed to ensure reliable and continuous operation under varying flow conditions. Operational considerations include maintaining consistent wastewater conveyance, minimizing system downtime, and allowing for routine inspection and maintenance. The incorporation of redundant pumping capacity and appropriate control of wet well levels helps ensure stable system performance and reduces the risk of service interruptions~\cite{TenStates_RecommendedStandards_2014,MDEQ_DesignGuidance_2025}.

Safety considerations focus on protecting both personnel and infrastructure during operation and maintenance activities. Pump stations are classified as confined spaces and may present hazards such as toxic gases, oxygen deficiency, and limited access. Therefore, the design must account for safe access, ventilation, and maintenance procedures in accordance with applicable safety regulations, including guidelines established by the Occupational Safety and Health Administration (OSHA) \cite{OSHA_ConfinedSpace_2024}. These considerations, along with adherence to wastewater design standards, help ensure safe operation and long-term system integrity~\cite{MDEQ_DesignGuidance_2025}.
